\documentclass[12pt]{article}
\usepackage[utf8]{inputenc}
\usepackage[T1]{fontenc}
\usepackage{hyperref}
\title{Marwari Language Resources for K-12 Curriculum Development}
\author{Compiled by Claude Sonnet 5.5 (claude-sonnet-5-5) for Andr\'as Kornai}
\date{}
\begin{document}
\maketitle

\section{Introduction}
Marwari (ISO 639-3 \texttt{rwr}) is an Indo-Aryan language of Rajasthan, India, written in Devanagari in India and Nastaliq in Pakistan. Wikipedia lists \texttt{rwr} under the broader Marwari code \texttt{mwr} and gives about 21 million native speakers (2011 census). The prompt's figure of about 25,000 differs by orders of magnitude, so the scope of \texttt{rwr} versus the macrolanguage needs checking. Dialects named include Thali/Bikaneri, Godwari, Dhatki and Shekhawati. I found no dedicated school textbooks in this search.

\section{Primary Educational Resources}
\begin{itemize}
\item Online Marwari lessons (speakmarwari.blogspot.co.uk) and a Marwari dictionary (hattai.page.tl) are linked from Omniglot; I did not open them. \url{https://omniglot.com/writing/marwari.htm}
\end{itemize}
I found no school textbooks or primers.

\section{Linguistic and Reference Materials}
\begin{itemize}
\item Wikipedia, Marwari language: \url{https://en.wikipedia.org/wiki/Marwari_language}
\item Omniglot, Marwari (writing and phrases, with links to external resources): \url{https://omniglot.com/writing/marwari.htm}
\item Universe of Memory, Marwari language resources. It lists an English grammar by Asopa (1950, on Scribd) and the eScholarship paper ``Topics in the Grammar of Marwari'': \url{https://universeofmemory.com/marwari-language-resources/}
\item Glosbe Marwari entry. The page I saw was titled as a Marwari--Likum dictionary, so check it before relying on it: \url{https://glosbe.com/rwr/lib}
\end{itemize}

\section{Religious and Community Resources}
None found in these results.

\section{Digital and Audio Resources}
\begin{itemize}
\item Wikitongues, Marwari (\texttt{rwr}) language page: \url{https://wikitongues.org/languages/rwr/}
\end{itemize}

\section{Curriculum Development Context}
The scope of the code \texttt{rwr} needs clarifying before planning. If it covers the large Marwari-speaking population of Rajasthan, the context is Hindi-medium schooling with Marwari mostly outside the curriculum. Existing materials found are reference, grammar and hobbyist lesson resources, with no school programme. A curriculum would have to be built from scratch around a chosen Devanagari orthography and a standard dialect.

\end{document}